\documentclass[twocolumn,aps,prx,10pt,floatfix]{revtex4-2}

\usepackage{graphicx}
\usepackage{amsmath,amssymb}
\usepackage{dcolumn}
\usepackage{bm}
\usepackage{hyperref}
\usepackage{xcolor}
\usepackage{microtype}

\graphicspath{
  {analysis_outputs/purification_total_charge_histogram_video_cpu/%
    N20_multi_geometry_nsh1_dwtrunc1_init-maxmix_S100_cycles-2Ny/figures/}
  {analysis_outputs/purification_total_charge_variance_histogram_video_cpu/%
    N20_multi_geometry_nsh1_dwtrunc1_init-maxmix_S100_cycles-2Ny/figures/}
  {analysis_outputs/purification_protocol_comparison/%
    N20_multi_geometry_nsh1_dwtrunc1_init-maxmix_S100_cycles-2Ny/figures/}
  {analysis_outputs/correlator_scaling_diagnostics/figures/}
  {analysis_outputs/perfect_correction_slope_excess_diagnostics/figures/}
  {analysis_outputs/purification_postselect_characterization/%
    N20_multi_geometry_nsh1_dwtrunc1_init-maxmix_S100_cycles-2Ny/figures/}
}

\begin{document}

\title{Charge Fluctuations in Class~A Fermionic Adaptive Circuits:\\
Purification Dynamics and Covariance-Matrix Scaling}

\author{A.~Bhuiyan}
\date{\today}

\begin{abstract}
We characterize charge fluctuations and entanglement structure in a two-dimensional
Class~A (U(1)-symmetric) adaptive Markov circuit acting on a free-fermion lattice.
Two complementary campaigns are analyzed.  In the first, the circuit is initialized
in the maximally mixed state and run under two distinct post-measurement protocols
(perfect correction and post-selection), revealing rapid suppression of charge
variance and entropy purification whose rates depend non-trivially on system
geometry.  In the second campaign, the circuit starts from a domain-wall initial
condition, and streaming covariance-matrix observables are recorded to extract the
scaling of the strip entanglement entropy and the off-diagonal square correlator
with subsystem size and inter-site separation, respectively.  The entropy obeys a
CFT-like chord-length formula with a fitted coefficient $c$, while the correlator
decays algebraically.  A persistent excess of the fitted correlator slope under
perfect correction relative to post-selection is characterized through a battery of
diagnostic tests.  Together these results probe the interplay of measurement-induced
criticality, topological structure, and charge conservation in a symmetry-protected
Gaussian circuit.
\end{abstract}

\maketitle

% ============================================================
\section{Introduction}
% ============================================================

Adaptive quantum circuits---in which mid-circuit measurements feed back adaptively
on subsequent unitary operations---have emerged as a rich arena for studying
measurement-induced phase transitions
(MIPTs)~\cite{Nahum2017,Li2018,Skinner2019,Chan2019,Bao2020}.
In free-fermion systems the Gaussian structure of the state is preserved
under Clifford/linear dynamics and projective measurements, enabling the
simulation of large system sizes that would otherwise be inaccessible.  When
a global U(1) charge symmetry is imposed---placing the circuit in \emph{Class~A}
of the ten-fold symmetry classification---the dynamics is further constrained:
the total fermion number is a conserved quantity at every step of the circuit,
and charge fluctuations become a direct probe of the structure of the
post-measurement ensemble.

In the present work we study a 2D rectangular lattice of size $N_x \times N_y$
with two fermionic orbitals (A and B) per unit cell, simulated via a Fermionic
Gaussian Tensor Network (FGTN) with nearest-neighbor Wannier support ($n_\text{sh}=1$).
The adaptive Markov cycle couples coherent evolution to local measurements; the
measurement outcomes can be handled in two limiting ways that bound the physical
ensemble:
\begin{itemize}
  \item \emph{Perfect correction}: measurement outcomes on every sample trajectory
    are post-processed so that the circuit corrects any unwanted charge offset,
    maintaining a neutral ensemble.  One hundred statistically independent samples
    are averaged.
  \item \emph{Post-selection}: a single trajectory is retained subject to
    post-selection on a designated measurement pattern.  The result is a single
    quantum trajectory sampled from the post-selected distribution.
\end{itemize}
These two protocols bracket the range of possible measurement-outcome averaging
and provide complementary windows onto the circuit dynamics.

The notes below are organized as follows.  In Sec.~\ref{sec:model} we define the
model and the covariance-matrix formalism.  Sections~\ref{sec:observables} and
\ref{sec:protocols} specify the observables and measurement protocols.
Section~\ref{sec:purification} presents the purification-dynamics campaign
(maxmix initialization) and Sec.~\ref{sec:streaming} presents the streaming
covariance campaign (domain-wall initialization).  Section~\ref{sec:diagnostics}
reports the slope-excess diagnostic analysis, and Sec.~\ref{sec:implications}
discusses the physical implications.

% ============================================================
\section{Model and Covariance-Matrix Formalism}
\label{sec:model}
% ============================================================

\subsection{Lattice and Hilbert space}

The system consists of $N_x \times N_y$ unit cells arranged on a rectangular
torus, with periodic boundary conditions in both directions.  Each unit cell
$r = (x,y)$ hosts two fermionic modes labelled by orbital indices $a$ and $b$,
giving a total of $N = 2 N_x N_y$ modes.  We denote the corresponding
creation/annihilation operators by $c_{r\mu}^\dagger$, $c_{r\mu}$ with
$\mu \in \{a,b\}$.  The global U(1) charge is
$\hat{Q} = \sum_{r,\mu} c_{r\mu}^\dagger c_{r\mu}$.

\subsection{Gaussian states and the covariance matrix}

All states encountered in the simulation are Gaussian (free-fermion) states,
completely characterized by the single-particle covariance matrix
\begin{equation}
  G_{\alpha\beta} \;=\;
  \langle c_\alpha c_\beta^\dagger \rangle - \tfrac{1}{2}\delta_{\alpha\beta},
  \label{eq:G}
\end{equation}
where the indices $\alpha, \beta$ run over all $N$ modes.
$G$ is Hermitian and satisfies $-\tfrac{1}{2}I \leq G \leq \tfrac{1}{2}I$.
The corresponding occupation matrix is
\begin{equation}
  C \;=\; \tfrac{1}{2}(G + I),
  \label{eq:C}
\end{equation}
with $0 \leq C \leq I$ in the sense of operators.  The mean occupation of mode
$\alpha$ is $\langle c_\alpha^\dagger c_\alpha\rangle = C_{\alpha\alpha}$.

The two limiting states used as initial conditions are:
\begin{itemize}
  \item \emph{Maximally mixed} ($G = 0$, $C = I/2$): every mode is
    independently at half-filling, with maximal single-mode entropy $\ln 2$.
  \item \emph{Default (domain-wall)} initial state: the circuit is initialized
    in a topological domain-wall configuration with domain-wall columns at
    $x_\text{DW} \in \{5, 15\}$, separating a topologically nontrivial interior
    from trivial edge regions.
\end{itemize}

The Markov cycle is implemented through the canonical entry point
\texttt{classA\_U1FGTN\_gpu.run\_markov\_circuit}, which performs a complete
sequence of coherent Wannier-function hops and adaptive single-site measurements
with $n_\text{sh} = 1$ nearest-neighbor support and domain-wall truncation.
Parameters $\alpha_1$ and $\alpha_2$ govern the measurement rate and adaptive
feedback strength.  GPU simulations ran on an NVIDIA A100 (40\,GiB) with batch
sizes tuned automatically to saturate memory.

% ============================================================
\section{Observables}
\label{sec:observables}
% ============================================================

\subsection{von Neumann entropy}

The total von Neumann entropy of a Gaussian state is
\begin{equation}
  S_\text{tot} \;=\;
  -\operatorname{Tr}\!\bigl[C \ln C + (I - C)\ln(I - C)\bigr],
  \label{eq:entropy}
\end{equation}
where the eigenvalues of $C$ lie in $[0,1]$ and are identified with
single-mode occupation numbers.  For the maxmix state $C = I/2$, yielding
$S_\text{tot}^\text{(mm)} = N\ln 2$.  For a pure state, all eigenvalues are
0 or 1 and $S_\text{tot} = 0$.

\subsection{Real-space Chern number}

The topological character of the occupied subspace is diagnosed by the
Kitaev real-space Chern number~\cite{Kitaev2006}
\begin{equation}
  \mathcal{C}_\text{RS} \;=\;
  12\pi i\,\operatorname{Tr}\!\bigl[
    \bar{C}_{CA}\bar{C}_{AB}\bar{C}_{BC}
    - \bar{C}_{AC}\bar{C}_{CB}\bar{C}_{BA}
  \bigr],
  \label{eq:chern}
\end{equation}
where $\bar{C} = C^\dagger$ and the subscript $XY$ restricts $\bar{C}$ to
rows in region $X$ and columns in region $Y$.  The three regions $A$, $B$, $C$
partition a disk of radius $r = 0.4\min(N_x,N_y)$ centered at the midpoint of
the system into three $120^\circ$ angular sectors.  For a topologically trivial
(Chern-0) insulator $\mathcal{C}_\text{RS} = 0$; for a $|\mathcal{C}|=1$ quantum
Hall state it returns $\pm1$.

\subsection{Local charge density and variance}

The mean charge in unit cell $r = (x,y)$ is
\begin{equation}
  \langle N_r \rangle \;=\; C_{a_r a_r} + C_{b_r b_r},
  \label{eq:localcharge}
\end{equation}
i.e., the sum of orbital-diagonal occupancies.  At half-filling
$\langle N_r \rangle = 1$ for all $r$.

The charge variance within a single unit cell is
\begin{equation}
  \operatorname{Var}(N_r)
  = \langle N_r\rangle + 2\!\left(C_{a_r a_r}C_{b_r b_r}
    - |C_{a_r b_r}|^2\right)
    - \langle N_r\rangle^2,
  \label{eq:localvar}
\end{equation}
which can be rewritten as
$C_{a_r a_r}(1-C_{a_r a_r}) + C_{b_r b_r}(1-C_{b_r b_r})
- 2|C_{a_r b_r}|^2$,
manifesting the Pauli-blocking reduction due to inter-orbital coherence.
For the maxmix state $C_{a_r b_r} = 0$ and $C_{a_r a_r} = C_{b_r b_r} = 1/2$,
giving $\operatorname{Var}(N_r)^\text{(mm)} = 1/2$.

\subsection{Total charge variance}

The global charge variance for a Gaussian state is
\begin{equation}
  \operatorname{Var}(\hat{Q}) \;=\;
  \operatorname{Tr}[C] - \operatorname{Tr}[C^2]
  \;=\; \operatorname{Tr}[C(I-C)],
  \label{eq:totalvar}
\end{equation}
which follows from the Wick theorem for free fermions.
The maxmix baseline is
$\operatorname{Var}(\hat{Q})^\text{(mm)} = N/4 = N_x N_y / 2$.
A fully charge-ordered (Fock) state satisfies $C^2 = C$, hence
$\operatorname{Var}(\hat{Q}) = 0$.

\subsection{Strip entanglement entropy and chord scaling}

For the streaming covariance campaign, the entanglement entropy of a horizontal
strip of width $A_y$ cells, with the strip bottom edge at $y_0$, is computed from
the restricted covariance block $G|_\text{strip}$.  Averaging over all $N_y$
reference positions,
\begin{equation}
  \overline{S}(A_y) \;=\; \frac{1}{N_y}\sum_{y_0=0}^{N_y-1} S(A_y; y_0).
  \label{eq:stripentropy}
\end{equation}
At a conformally invariant critical point in effectively (1+1)-dimensional
imaginary-time evolution, the entropy scales as~\cite{Calabrese2004}
\begin{equation}
  \overline{S}(A_y) \;\approx\; c \ln\!\left[
    \frac{N_y}{\pi}\sin\!\frac{\pi A_y}{N_y}\right] + \text{const},
  \label{eq:cftscaling}
\end{equation}
where the prefactor $c$ is related to the central charge of the effective CFT.
We extract $c$ by fitting $\overline{S}$ against the log-chord variable
$\xi = \ln\!\left[(N_y/\pi)\sin(\pi A_y/N_y)\right]$ over the window
$A_y \in [8, N_y/2]$.

\subsection{Off-diagonal square correlator}

A second probe of correlations is the $x$-averaged off-diagonal square
correlator at $y$-separation $r_y$:
\begin{equation}
  \widetilde{C}^2(r_y) \;=\;
  \frac{1}{2N_xN_y}
  \sum_{x,y,\mu,\nu}
  \bigl|C_{\mu(x,y),\,\nu(x,\,y+r_y)}\bigr|^2,
  \label{eq:correlator}
\end{equation}
where the orbital indices $\mu,\nu \in \{a,b\}$ and the site index
$(x,y+r_y)$ is taken modulo $N_y$.  In a critical phase one expects
algebraic decay
\begin{equation}
  \widetilde{C}^2(r_y) \;\propto\; r_y^{\alpha},\quad \alpha < 0,
  \label{eq:powerlaw}
\end{equation}
and we extract $\alpha$ by a log-log linear fit over
$r_y \in [5, N_y/2]$.

% ============================================================
\section{Measurement Protocols}
\label{sec:protocols}
% ============================================================

At the end of each Markov cycle, the Born-rule probabilities for the
adaptive measurements are recorded as the four maps
$p_\text{occ}(A^-_{r}|G)$,
$p_\text{occ}(B^-_{r}|G)$,
$1 - p_\text{occ}(A^+_{r}|G)$,
$1 - p_\text{occ}(B^+_{r}|G)$,
labelled $N^-_A$, $N^-_B$, $1-N^+_A$, $1-N^+_B$ in the data.
These quantify how strongly each site is being measured toward the occupied
or empty branch and serve as indicators of measurement frustration.

\paragraph{Perfect correction.}
The circuit evolves 100 statistically independent samples in parallel
(batched GPU execution).  At each cycle, the Markov update is applied and
the resulting covariance matrix is observed; any charge offset produced by
measurement is corrected by a classically controlled feedback unitary.
Sample-averaged quantities $\langle \mathcal{O}\rangle_S$ are reported.
This protocol is the physically relevant ensemble average and corresponds
to the mixed state $\rho = S^{-1}\sum_s |\psi_s\rangle\langle\psi_s|$.

\paragraph{Post-selection.}
A single trajectory is propagated with post-selection: at each adaptive
measurement site, the lower-energy (or fixed-pattern) branch is selected.
The result is one quantum trajectory drawn from the strongly post-selected
distribution, analogous to a ground-state projection in imaginary time.
Only one sample is generated since post-selection renders the trajectory
unique given the initial condition and pattern.

The two protocols are systematically compared by overlaying their
time-series and final-cycle spatial maps on the same axes.

% ============================================================
\section{Campaign~I: Purification from the Maximally Mixed State}
\label{sec:purification}
% ============================================================

\subsection{Setup}

The first campaign\footnote{%
  Campaign id:
  \texttt{N20\_multi\_geometry\_nsh1\_dwtrunc1%
  \_init-maxmix\_S100\_cycles-2Ny}.}
initializes all samples in the maxmix state $G = 0$ and evolves under the
Markov cycle for $T = 2N_y$ cycles.  Three geometries are studied:
$N_x = 20$ with $N_y \in \{30, 40, 50\}$.  Perfect correction runs use
$S = 100$ samples; the post-selection run uses a single trajectory.

The maxmix starting values of the diagnostics are:
\begin{align}
  S_\text{tot}^\text{(mm)} &= N\ln 2 = 2N_xN_y\ln 2,\\
  \operatorname{Var}(\hat{Q})^\text{(mm)} &= N_xN_y/2.
\end{align}
For $N_x = 20$, $N_y = 40$ these are approximately 1109~nats and 400,
respectively.

\subsection{Entropy purification}

\begin{figure}[!ht]
  \centering
  \includegraphics[width=\columnwidth]{%
    purification_total_entropy_sample_mean_vs_cycle_loglog.png}
  \caption{%
    Sample-averaged total von Neumann entropy
    $\langle S_\text{tot}\rangle_S$ vs.\ Markov cycle on a log-log scale,
    for all three geometries under perfect correction.
    The entropy falls precipitously from the maxmix value within the
    first few cycles and plateaus at an exponentially small residual.
  }
  \label{fig:entropy}
\end{figure}

Figure~\ref{fig:entropy} shows the sample-averaged entropy as a function
of Markov cycle on a log-log scale.
The entropy collapses by roughly an order of magnitude per cycle in the
earliest cycles, reaching $\langle S_\text{tot}\rangle \approx 0.3$~nats
for $N_x=20$, $N_y=40$ after $T=80$ cycles (compared to the maxmix value
of $\sim 1109$~nats).  This indicates that the Markov circuit is an
extremely efficient purifier: the adaptive measurements project the
maxmix density matrix essentially to a pure (or near-pure) state within
$O(N_y)$ cycles.

The $N_y$ dependence of the saturation value is mild; larger $N_y$ systems
retain slightly more residual entropy at the same cycle count, consistent
with the intuition that longer systems have more modes to disentangle.

\subsection{Charge-variance suppression}

\begin{figure}[!ht]
  \centering
  \includegraphics[width=\columnwidth]{%
    purification_total_charge_variance_raw_mean_vs_cycle.png}
  \caption{%
    Sample-averaged total charge variance
    $\langle\operatorname{Var}(\hat{Q})\rangle_S$ vs.\ Markov cycle for
    the three geometries under perfect correction.
    The variance drops from the maxmix baseline $N_xN_y/2$ to
    $O(10^{-2})$ within $2N_y$ cycles, demonstrating rapid charge
    crystallization.
  }
  \label{fig:chargevar}
\end{figure}

The total charge variance (Fig.~\ref{fig:chargevar}) follows a similar
rapid purification.  Starting from $\operatorname{Var}(\hat{Q})^\text{(mm)}
= N_xN_y/2$ (400 for $N_y=40$), the variance is already reduced to $\sim 41$
after the first cycle and reaches $\approx 0.08$ at the final cycle---a
five-thousand-fold suppression.  This is consistent with the circuit driving
the system into a near-Fock state in which each site has a definite charge.

The percentage charge variance $\operatorname{Var}(\hat{Q}) / (N_xN_y/2)$
shows that all three $N_y$ geometries reach comparably small residual fractions,
with the approach becoming slightly slower for larger $N_y$.

\subsection{Sample-to-sample fluctuations of the total charge}

Under perfect correction, the circuit does not enforce a fixed total charge;
rather, it corrects the measurement back-action on average.  The centered
total charge is defined as
\begin{equation}
  Q_\text{raw} \;=\; 2\!\left(\sum_r \langle N_r\rangle - N_xN_y\right),
  \label{eq:centeredcharge}
\end{equation}
which vanishes at half-filling, and its percentage counterpart
$Q_\text{pct} = 100 \cdot Q_\text{raw}/(2N_xN_y)$ measures the deviation
from half-filling as a fraction of the total mode count.
A nonzero sample-to-sample variance in $Q_\text{raw}$ indicates that
individual trajectories may carry a net charge offset even though the
ensemble mean is controlled.

From the data, the sample mean of $Q_\text{raw}$ tracks to zero throughout
the circuit evolution (the perfect-correction protocol enforces approximate
charge neutrality by construction), while the sample-to-sample variance
decreases monotonically, indicating that the perfect-correction feedback is
effective at collimating the charge distribution toward zero offset.

\subsection{Protocol comparison: perfect correction vs.\ post-selection}

\begin{figure}[!ht]
  \centering
  \includegraphics[width=\columnwidth]{%
    purification_protocol_comparison_entropy_loglog.png}
  \caption{%
    Log-log comparison of total entropy vs.\ cycle for perfect correction
    (100-sample average, shown with $\pm1\sigma$ band) and post-selection
    (single trajectory) for all three geometries.
    The post-selected trajectory purifies more rapidly and to a smaller
    residual entropy at early cycles, reflecting the more aggressive
    projection of the post-selection protocol.
  }
  \label{fig:protocolentropy}
\end{figure}

\begin{figure}[!ht]
  \centering
  \includegraphics[width=\columnwidth]{%
    purification_protocol_comparison_charge_variance_loglog.png}
  \caption{%
    Log-log comparison of total charge variance vs.\ cycle for the two
    protocols.  The post-selected trajectory again purifies faster,
    reaching smaller residual variance than the 100-sample perfect-correction
    average.
  }
  \label{fig:protocolvar}
\end{figure}

Figures~\ref{fig:protocolentropy} and \ref{fig:protocolvar} directly
compare the two protocols.  The post-selected trajectory consistently
purifies faster and to a smaller residual than the perfect-correction
ensemble average.  This is physically expected: post-selection biases toward
trajectories that happen to undergo rapid purification, whereas the
perfect-correction ensemble contains trajectories with varying purification
rates.

\begin{figure}[!ht]
  \centering
  \includegraphics[width=\columnwidth]{%
    purification_protocol_comparison_chern_loglog.png}
  \caption{%
    Log-log comparison of the real-space Chern number
    $\mathcal{C}_\text{RS}$ vs.\ cycle for the two protocols.
    Both converge to an integer value (consistent with $|\mathcal{C}|=1$)
    at late cycles, confirming that the steady state is topologically
    nontrivial.
  }
  \label{fig:protocolchern}
\end{figure}

Figure~\ref{fig:protocolchern} shows the real-space Chern number for both
protocols.  At late cycles both converge to a value consistent with
$|\mathcal{C}_\text{RS}| = 1$, indicating that the adaptive circuit
drives the system toward a topological (Chern-insulator-like) steady state.
The Chern number build-up is faster for post-selection, consistent with the
faster purification dynamics observed in the entropy and charge variance.

\subsection{Spatial charge density and variance}

The sample-averaged local charge density
$\langle N_r\rangle = \frac{1}{S}\sum_s C^{(s)}_{a_r a_r} + C^{(s)}_{b_r b_r}$
evolves from the uniform half-filling $\langle N_r\rangle = 1$ of the
maxmix state and, at late times, develops spatial structure reflecting the
domain-wall geometry encoded in the adaptive feedback.  Heatmap videos
of $\langle N_r\rangle$ over cycles confirm that the charge density
settles into a spatially non-uniform pattern commensurate with the
nontrivial winding of the Chern insulator state.

Similarly, the local charge variance map
$\operatorname{Var}(N_r)$ (cf.\ Eq.~\eqref{eq:localvar}) is initially $1/2$
everywhere and collapses to near zero at late times, with small residual
hot spots at the domain-wall columns $x \in \{5, 15\}$, consistent with
the expectation that domain-wall modes (edge states) retain larger charge
fluctuations than bulk modes.

\clearpage
% ============================================================
\section{Campaign~II: Streaming Covariance Observables}
\label{sec:streaming}
% ============================================================

\subsection{Setup}

The second campaign\footnote{%
  Campaign id:
  \texttt{N20\_selected\_nsh1\_dwtrunc1%
  \_init-default\_S100\_cycles-2Ny}.}
uses the domain-wall initial condition with $\mathbf{x}_\text{DW} = (5, 15)$
and runs for $T = 2N_y$ Markov cycles.  The three complete geometries are
$N_y \in \{30, 40, 50\}$ (with $N_x = 20$ throughout); a fourth run with
$N_y = 60$ was initiated but did not complete.  One hundred samples are
accumulated for the perfect-correction protocol; a single post-selected
trajectory is also recorded.

Rather than saving the full covariance matrix history, this campaign uses
a \emph{streaming} observer that computes and accumulates the following
derived quantities at each cycle without retaining $G(t)$:
the strip entanglement entropy $\overline{S}(A_y)$, the off-diagonal square
correlator $\widetilde{C}^2(r_y)$, the local charge-cell occupancy map, and
the measurement frustration.

\subsection{Strip entropy and CFT scaling}

\begin{figure}[!ht]
  \centering
  \includegraphics[width=\columnwidth]{%
    equilibrium_vs_dynamics_correlator_curves.png}
  \caption{%
    Sample-averaged off-diagonal square correlator
    $\widetilde{C}^2(r_y)$ vs.\ $r_y$ at the final cycle
    ($t = 2N_y$) for the three geometries under perfect correction.
    On a log-log scale the curves are approximately linear,
    supporting the power-law ansatz $\widetilde{C}^2 \propto r_y^\alpha$.
  }
  \label{fig:correlators}
\end{figure}

\begin{figure}[!ht]
  \centering
  \includegraphics[width=\columnwidth]{%
    n20_default_protocol_comparison.png}
  \caption{%
    Comparison of the fitted correlator log-log slope $\alpha$ (top panel)
    and the entropy chord-scaling coefficient $c$ (bottom panel) for
    perfect correction and post-selection as functions of cycle.
    The perfect-correction slope is systematically more negative (steeper
    decay) than post-selection throughout the evolution.
  }
  \label{fig:protocolcomparison}
\end{figure}

The strip entropy data are fitted cycle-by-cycle to the chord formula
Eq.~\eqref{eq:cftscaling} over the window $A_y \in [8, N_y/2]$.
The extracted coefficient $c(t)$ grows from a small initial value and
saturates at late cycles; its $N_y$ dependence distinguishes between a
volume-law ($c \propto N_y$), critical ($c \sim$ const), or area-law
($c \to 0$) phase.  Preliminary results (Fig.~\ref{fig:protocolcomparison},
lower panel) indicate that $c$ stabilizes at a finite $O(1)$ value
for the geometries studied here, suggestive of a critical or weakly
volume-law phase.

\subsection{Off-diagonal square correlator and power-law exponent}

The off-diagonal square correlator (Eq.~\eqref{eq:correlator}) is shown
in Fig.~\ref{fig:correlators} at the final cycle for the three $N_y$
geometries.  On a log-log scale the curves are approximately linear
over the fitting window $r_y \in [5, N_y/2]$, supporting the power-law
ansatz Eq.~\eqref{eq:powerlaw}.

The fitted exponent $\alpha$ evolves with cycle count (upper panel of
Fig.~\ref{fig:protocolcomparison}).  For the perfect-correction protocol,
$\alpha$ is consistently more negative (faster algebraic decay) than for
post-selection.  This slope excess is a non-trivial feature of the ensemble
averaging and is investigated in detail in Sec.~\ref{sec:diagnostics}.

\subsection{Fit-window sensitivity}

\begin{figure}[!ht]
  \centering
  \includegraphics[width=\columnwidth]{n20_final_cycle_window_heatmap.png}
  \caption{%
    Heatmap of the fitted correlator slope $\alpha$ as a function of
    the minimum fitting window boundary $r_y^\text{min}$ (vertical axis)
    and Markov cycle (horizontal axis) for the $N_x=20$, $N_y=40$ perfect-correction
    run.  The slope is relatively stable for $r_y^\text{min} \lesssim 10$
    but shows increased scatter near the ultraviolet cutoff.
  }
  \label{fig:windowheatmap}
\end{figure}

To assess the robustness of the extracted exponent, the fit is repeated
for varying minimum window boundaries $r_y^\text{min}$.  The resulting
slope heatmap (Fig.~\ref{fig:windowheatmap}) shows that the exponent is
stable over a wide range of window choices at late cycles, while early-cycle
transients produce more scatter.  This validates the default fitting window
and confirms that the late-cycle exponent is not an artifact of the
specific choice of $r_y^\text{min}$.

\clearpage
% ============================================================
\section{Slope-Excess Diagnostics}
\label{sec:diagnostics}
% ============================================================

\subsection{Observation}

The systematic excess of the perfect-correction correlator slope
$|\alpha_\text{pc}| > |\alpha_\text{ps}|$ observed in
Fig.~\ref{fig:protocolcomparison} raises the question of whether it is a
genuine physical effect or an artifact of the ensemble construction.  Possible
sources include:
(i) sample-to-sample correlations introduced by the perfect-correction
feedback, which may bias the averaged $C^2$ relative to the true mixed
state correlator;
(ii) charge-sector contamination from trajectories with nonzero net charge;
(iii) a difference in how the two protocols explore the measurement-outcome
distribution.

\subsection{Diagnostic tests}

\begin{figure}[!ht]
  \centering
  \includegraphics[width=\columnwidth]{test3_delta_slope_vs_cycle.png}
  \caption{%
    Delta slope $\Delta\alpha = \alpha_\text{pc} - \alpha_\text{ps}$ as a
    function of Markov cycle for the three $N_y$ geometries.  The excess
    grows from zero at early cycles and stabilizes at a negative value at
    late cycles, indicating that the slope difference is a genuine feature
    of the late-time dynamics rather than an early-time transient.
  }
  \label{fig:deltaslope}
\end{figure}

\begin{figure}[!ht]
  \centering
  \includegraphics[width=\columnwidth]{test4_averaging_order.png}
  \caption{%
    Comparison of the fitted slope obtained by (a) fitting each sample
    individually and then averaging the slopes, vs.\ (b) averaging the
    correlator curves and then fitting the average.  The two procedures
    give consistent results, ruling out Jensen's inequality as the source
    of the slope excess.
  }
  \label{fig:avgorder}
\end{figure}

We perform five diagnostic tests:

\paragraph{Test~1 (metric correlations).}
The excess slope $\Delta\alpha$ is regressed against scalar metrics
(entropy, Chern number, total charge variance) across the 100 samples.
No single metric explains the bulk of the variance in $\Delta\alpha$,
suggesting the excess is not driven by a simple correlation with
purification speed.

\paragraph{Test~2 (conditioned subsets).}
The 100 perfect-correction samples are stratified by entropy quartile
at the final cycle; the correlator slope is fitted separately in each
quartile.  The slope varies weakly across quartiles, indicating that
the excess is not exclusively contributed by slowly purifying outliers.

\paragraph{Test~3 (cycle dependence).}
Figure~\ref{fig:deltaslope} shows $\Delta\alpha(t) = \alpha_\text{pc}(t)
- \alpha_\text{ps}(t)$ as a function of Markov cycle.  The excess grows
from zero in the early cycles and saturates at a negative value
($|\alpha_\text{pc}| > |\alpha_\text{ps}|$) by mid-circuit, persisting
to the final cycle.  This rules out an early-time initialization artifact.

\paragraph{Test~4 (averaging order).}
Figure~\ref{fig:avgorder} compares two procedures: fitting each sample
then averaging the slopes (``slope average'') vs.\ averaging the
correlator curves then fitting (``average then fit'').  The two give
consistent values, ruling out a Jensen-inequality bias as the dominant
source of the excess.

\paragraph{Test~5 (neutralized ensemble).}
A subset of samples with near-zero centered total charge ($|Q_\text{pct}|
< \epsilon$) is selected, and the correlator slope is refit.  The slope
excess persists in this neutralized ensemble, indicating that charge-sector
contamination from off-half-filling trajectories is not the primary driver.

\subsection{Interpretation}

The diagnostic battery collectively suggests that the slope excess is a
genuine property of the perfect-correction protocol ensemble rather than a
numerical or sampling artifact.  A plausible physical mechanism is the
following: under perfect correction, the feedback unitary applied at each
cycle introduces a coherent phase shift that is \emph{correlated} across
the sample batch---the correction is derived from the same measurement
outcomes that determined the post-measurement state.  This correlation
enriches the off-diagonal structure of the ensemble-averaged occupation
matrix relative to the uncorrelated post-selected trajectory, producing a
faster algebraic decay of $\widetilde{C}^2(r_y)$.

An alternative explanation involves the boundary conditions on the charge
sector: perfect correction enforces approximate $\hat{Q} \approx N_xN_y$
at the ensemble level, which may impose additional off-diagonal coherences
across unit cells that the post-selected trajectory (which selects one
charge sector deterministically) does not share.  Distinguishing these
mechanisms conclusively would require access to the full sample covariance
matrix $\langle G_{\alpha\beta} G_{\gamma\delta}\rangle_S$, which is not
currently saved.

% ============================================================
\section{Physical Implications}
\label{sec:implications}
% ============================================================

\paragraph{Rapid purification and topological ordering.}
The most striking observation from Campaign~I is the speed of purification.
Within $O(N_y)$ adaptive Markov cycles the system transitions from a
maximally disordered maxmix state to a near-pure Chern-insulator state
with $|\mathcal{C}_\text{RS}|\approx 1$ and charge variance
$\ll N_xN_y/2$.  This is consistent with the theoretical expectation that
adaptive circuits in the ``area-law'' phase of the MIPT can purify
exponentially fast~\cite{Gullans2020}.  The convergence of the real-space
Chern number confirms that the steady state selected by the adaptive
feedback is topologically nontrivial---the measurement back-action drives
the system toward a specific topological phase rather than a featureless
Fock state.

\paragraph{Charge conservation and measurement-induced order.}
The U(1) charge symmetry of the Class~A circuit has a direct effect on
the charge-variance dynamics.  The rapid suppression of $\operatorname{Var}
(\hat{Q})$ from the maxmix baseline is not merely purification (which would
decrease all entropy sources equally) but reflects a specific \emph{charge
crystallization}: individual measurement trajectories localize in a narrow
charge sector much more rapidly than a symmetry-breaking circuit would.
The persistence of small residual charge variance ($\sim 10^{-2}$) at late
times under perfect correction, compared to the even smaller residual under
post-selection, quantifies how strongly the ensemble spread in charge
deviates from the post-selected pure-state value.

\paragraph{Critical correlations in the domain-wall geometry.}
Campaign~II reveals that in the presence of a topological domain wall,
the circuit generates long-range algebraic correlations in the off-diagonal
covariance $\widetilde{C}^2(r_y) \propto r_y^\alpha$.  The fact that the
entropy grows logarithmically with chord length (consistent with
Eq.~\eqref{eq:cftscaling} with finite $c$) and the correlator decays
as a power law are hallmarks of a critical or weakly volume-law entangled
state.  Whether this represents true criticality (i.e., the circuit sits
at or near a measurement-induced critical point) or a crossover between
a bulk topological phase and edge-localized modes is not yet determined;
a systematic $N_y$ scaling analysis of $c$ and $\alpha$ is needed to
distinguish these scenarios.

\paragraph{Protocol dependence and its diagnostic implications.}
The systematic difference in correlator scaling between perfect correction
and post-selection (Sec.~\ref{sec:diagnostics}) has implications for the
interpretation of circuit simulations.  If the slope excess is physical
rather than artifactual, it implies that the measurement-averaged ensemble
(perfect correction) and the individual trajectory (post-selection) belong
to different universality classes of the MIPT, in agreement with the
general principle that the MIPT is a property of the distribution
of quantum states rather than any single trajectory.  Verifying this
interpretation requires larger system sizes and, ideally, access to the
replica structure~\cite{Bao2020,Jian2020} of the measurement-averaged state.

\paragraph{Future directions.}
Several extensions are immediately motivated:
(i) completing the $N_y = 60$ streaming covariance run to extend the
finite-size scaling of $c$ and $\alpha$;
(ii) computing the two-copy (second-R\'enyi) entropy to access the
entanglement negativity spectrum and probe the MIPT universality class more
directly;
(iii) analyzing the spatial charge density profile near the domain wall to
extract the localization length of the edge modes as a function of cycle;
(iv) performing a finite-size scaling collapse of $\widetilde{C}^2(r_y)$
to extract the critical exponent $\nu$ of the MIPT.

% ============================================================
\begin{acknowledgments}
Computations were performed on NVIDIA A100 GPUs via Google Colab.
\end{acknowledgments}

\begin{thebibliography}{99}

\bibitem{Nahum2017}
A.~Nahum, J.~Ruhman, S.~Vijay, and J.~Haah,
\textit{Quantum Entanglement Growth Under Random Unitary Dynamics},
Phys.\ Rev.\ X \textbf{7}, 031016 (2017).

\bibitem{Li2018}
Y.~Li, X.~Chen, and M.~P.~A.~Fisher,
\textit{Quantum Zeno effect and the many-body entanglement transition},
Phys.\ Rev.\ B \textbf{98}, 205136 (2018).

\bibitem{Skinner2019}
B.~Skinner, J.~Ruhman, and A.~Nahum,
\textit{Measurement-Induced Phase Transitions in the Dynamics of Entanglement},
Phys.\ Rev.\ X \textbf{9}, 031009 (2019).

\bibitem{Chan2019}
A.~Chan, R.~M.~Nandkishore, M.~Pretko, and G.~Smith,
\textit{Unitary-projective entanglement dynamics},
Phys.\ Rev.\ B \textbf{99}, 224307 (2019).

\bibitem{Bao2020}
Y.~Bao, S.~Choi, and E.~Altman,
\textit{Theory of the phase transition in random unitary circuits with measurements},
Phys.\ Rev.\ B \textbf{101}, 104301 (2020).

\bibitem{Kitaev2006}
A.~Kitaev,
\textit{Anyons in an exactly solved model and beyond},
Ann.\ Phys.\ \textbf{321}, 2 (2006).

\bibitem{Calabrese2004}
P.~Calabrese and J.~Cardy,
\textit{Entanglement entropy and quantum field theory},
J.\ Stat.\ Mech.\ \textbf{2004}, P06002 (2004).

\bibitem{Gullans2020}
M.~J.~Gullans and D.~A.~Huse,
\textit{Dynamical Purification Phase Transition Induced by Quantum Measurements},
Phys.\ Rev.\ X \textbf{10}, 041020 (2020).

\bibitem{Jian2020}
C.-M.~Jian, Y.-Z.~You, R.~Vasseur, and A.~W.~W.~Ludwig,
\textit{Measurement-induced criticality in random quantum circuits},
Phys.\ Rev.\ B \textbf{101}, 104302 (2020).

\end{thebibliography}

\end{document}
